\documentclass[10pt,a4paper]{amsart}
\usepackage{amsmath,amssymb,amsthm,graphicx}
% For commutative diagrams we use Xy-pic
\usepackage[all]{xy}

% We use 2cell for 2-commutative diagrams.
\xyoption{2cell}
% Theorem environments.
%
\theoremstyle{plain}
\newtheorem{theorem}[subsection]{Theorem}
\newtheorem{proposition}[subsection]{Proposition}
\newtheorem{lemma}[subsection]{Lemma}

\theoremstyle{definition}
\newtheorem{definition}[subsection]{Definition}
\newtheorem{example}[subsection]{Example}
\newtheorem{exercise}[subsection]{Exercise}
\newtheorem{situation}[subsection]{Situation}

\theoremstyle{remark}
\newtheorem{remark}[subsection]{Remark}
\newtheorem{remarks}[subsection]{Remarks}

\numberwithin{equation}{subsection}

% Macros
%
\def\lim{\mathop{\mathrm{lim}}\nolimits}
\def\colim{\mathop{\mathrm{colim}}\nolimits}
\def\Spec{\mathop{\mathrm{Spec}}}
\def\Hom{\mathop{\mathrm{Hom}}\nolimits}
\def\Ext{\mathop{\mathrm{Ext}}\nolimits}
\def\SheafHom{\mathop{\mathcal{H}\!\mathit{om}}\nolimits}
\def\SheafExt{\mathop{\mathcal{E}\!\mathit{xt}}\nolimits}
\def\Sch{\mathit{Sch}}
\def\Mor{\mathop{\mathrm{Mor}}\nolimits}
\def\Ob{\mathop{\mathrm{Ob}}\nolimits}
\def\Sh{\mathop{\mathit{Sh}}\nolimits}
\def\NL{\mathop{N\!L}\nolimits}
\def\CH{\mathop{\mathrm{CH}}\nolimits}
\def\proetale{{pro\text{-}\acute{e}tale}}
\def\etale{{\acute{e}tale}}
\def\QCoh{\mathit{QCoh}}
\def\Ker{\mathop{\mathrm{Ker}}}
\def\Im{\mathop{\mathrm{Im}}}
\def\Coker{\mathop{\mathrm{Coker}}}
\def\Coim{\mathop{\mathrm{Coim}}}

% Boxtimes
%
\DeclareMathSymbol{\boxtimes}{\mathbin}{AMSa}{"02}

%
% Macros for moduli stacks/spaces
%
\def\QCohstack{\mathcal{QC}\!\mathit{oh}}
\def\Cohstack{\mathcal{C}\!\mathit{oh}}
\def\Spacesstack{\mathcal{S}\!\mathit{paces}}
\def\Quotfunctor{\mathrm{Quot}}
\def\Hilbfunctor{\mathrm{Hilb}}
\def\Curvesstack{\mathcal{C}\!\mathit{urves}}
\def\Polarizedstack{\mathcal{P}\!\mathit{olarized}}
\def\Complexesstack{\mathcal{C}\!\mathit{omplexes}}
% \Pic is the operator that assigns to X its picard group, usage \Pic(X)
% \Picardstack_{X/B} denotes the Picard stack of X over B
% \Picardfunctor_{X/B} denotes the Picard functor of X over B
\def\Pic{\mathop{\mathrm{Pic}}\nolimits}
\def\Picardstack{\mathcal{P}\!\mathit{ic}}
\def\Picardfunctor{\mathrm{Pic}}
\def\Deformationcategory{\mathcal{D}\!\mathit{ef}}
\usepackage{cleveref}
\crefname{originallemma}{original lemma}{original lemmas}
\Crefname{originallemma}{Original lemma}{Original lemmas}
\makeatletter
\newlabel{algebra-lemma-base-change-strict-henselization-quasi-finite}{{0GIP}{source}{}{sourcepointer-1}{}}
\newlabel{algebra-lemma-base-change-strict-henselization-quasi-finite@cref}{{[originallemma][0][]0GIP}{[1][0][]source}{}{}{}}
\makeatother
\title{Local acyclicity is stable under a base change that is quasi-finite, an affine inverse limit of quasi-finite schemes, or integral}
\author{The Stacks Project Authors}
\date{stacks-project:etale-cohomology.tex:local-snapshot-20261011; GNU FDL 1.2 or later}
\begin{document}
\maketitle
\noindent\textit{Editorial note.} All delivered mathematical text is the author's own, copied from the pinned Stacks Project chapter file etale-cohomology.tex; no statement or proof step was rewritten, reordered or removed.\par\medskip
\noindent\textit{Editorial note.} The counted proof is the author's complete argument: the proof is written out in full: the comparison of strict henselisations under the three alternative hypotheses on g, the identification of the residue fields, and the resulting equality of the varieties of vanishing cycles.\par\medskip
\noindent\textit{Editorial note.} The label lemma-locally-acyclic-quasi-finite-base-change names the counted statement, and the proof environment immediately following it in the source is delivered verbatim as its proof.\par\medskip
\noindent\textit{Editorial note.} The counted statement is carried by the author's own lemma environment in the chapter source.\par\medskip
\noindent\textit{Editorial note.} The statement and proof are copied from the pinned chapter file; the single cited input is kept as an external pointer with its canonical Stacks tag.\par\medskip
\noindent\textit{Editorial note.} The preamble is the chapter's own native preamble excerpt already used for this source: the Xy-pic setup (native preamble lines 14-18) and the native theorem declarations with the mathematical macros (native preamble lines 154-217).\par\medskip
\noindent\textit{Allowed standard prerequisites (editorial).}\begin{itemize}
\item Base change of strict henselisations along a quasi-finite morphism. Algebra, the lemma identifying O\textasciicircum{}sh\_\{X,x\} tensor\_\{O\textasciicircum{}sh\_\{S,s\}\} O\textasciicircum{}sh\_\{S',s'\} with O\textasciicircum{}sh\_\{X',x'\}. Supplies the identification of the strict henselisations that is the technical heart of the proof.
\item Varieties of vanishing cycles and the definition of local acyclicity. Same chapter, the definition of local acyclicity at a geometric point relative to a complex K through the varieties F\_\{x,t\}. Supplies the objects whose identification gives the conclusion.
\end{itemize}\medskip
\ifdefined\hypertarget\hypertarget{sourcepointer-1}{}\fi
\noindent\textit{Source reference.} Original lemma 0GIP: sources/stacks-chow/tags.txt, source lines 20624--20624. Prerequisite pointer invoked at proof line(s) 19316; it names a standard fact of a sibling chapter that the counted argument uses, and every step of the counted proof built on it is retained verbatim in the delivered body.\par\medskip
\begin{lemma}
\label{lemma-locally-acyclic-quasi-finite-base-change}
Let
$$
\xymatrix{
X' \ar[r]_{g'} \ar[d]_{f'} & X \ar[d]^f \\
S' \ar[r]^g & S
}
$$
be a cartesian diagram of schemes. Let $K$ be an object of $D(X_\etale)$.
Let $\overline{x}'$ be a geometric point of $X'$ with image $\overline{x}$
in $X$. If
\begin{enumerate}
\item $f$ is locally acyclic at $\overline{x}$ relative to $K$ and
\item $g$ is locally quasi-finite, or $S' = \lim S_i$
is a directed inverse limit of schemes locally quasi-finite over $S$
with affine transition morphisms, or $g : S' \to S$ is integral,
\end{enumerate}
then $f'$ locally acyclic at $\overline{x}'$ relative to $(g')^{-1}K$.
\end{lemma}

\begin{proof}
Denote $\overline{s}'$ and $\overline{s}$ the images of
$\overline{x}'$ and $\overline{x}$ in $S'$ and $S$.
Let $\overline{t}'$ be a geometric point of the spectrum of
$\Spec(\mathcal{O}^{sh}_{S', \overline{s}'})$ and denote
$\overline{t}$ the image in $\Spec(\mathcal{O}^{sh}_{S, \overline{s}})$.
By Algebra, Lemma
\ref{algebra-lemma-base-change-strict-henselization-quasi-finite}
and our assumptions on $g$ we have
$$
\mathcal{O}^{sh}_{X, \overline{x}}
\otimes_{\mathcal{O}^{sh}_{S, \overline{s}}}
\mathcal{O}^{sh}_{S', \overline{s}'}
\longrightarrow
\mathcal{O}^{sh}_{X', \overline{x}'}
$$
is an isomorphism. Since by our conventions
$\kappa(\overline{t}) = \kappa(\overline{t}')$
we conclude that
$$
F_{\overline{x}', \overline{t}'} =
\Spec\left(
\mathcal{O}^{sh}_{X', \overline{x}'}
\otimes_{\mathcal{O}^{sh}_{S', \overline{s}'}}
\kappa(\overline{t}')\right) =
\Spec\left(
\mathcal{O}^{sh}_{X, \overline{x}}
\otimes_{\mathcal{O}^{sh}_{S, \overline{s}}}
\kappa(\overline{t})\right) =
F_{\overline{x}, \overline{t}}
$$
In other words, the varieties of vanishing cycles
of $f'$ at $\overline{x}'$ are examples of varieties of vanishing
cycles of $f$ at $\overline{x}$. The lemma follows
immediately from this and the definitions.
\end{proof}

\end{document}
